\documentclass[11pt]{article}
\usepackage[a4paper,margin=2.05cm]{geometry}
\usepackage[T1]{fontenc}
\usepackage{lmodern}
\usepackage{amsmath,amssymb,amsthm}
\usepackage{xcolor}
\usepackage{microtype}
\usepackage[hidelinks]{hyperref}
\setlength{\emergencystretch}{3em}

\definecolor{ink}{HTML}{172033}
\definecolor{accent}{HTML}{385B7A}
\definecolor{soft}{HTML}{EEF5FA}
\color{ink}

\newtheorem*{theorem}{Literature-implied theorem}
\newcommand{\E}{\mathbb E}

\title{The full power growth of $k$-wise Khinchine constants}
\author{Literature-implied answer to the questions in arXiv:1412.7859}
\date{11 August 2026}

\begin{document}
\maketitle

\begin{center}
\fcolorbox{accent}{soft}{\parbox{0.91\textwidth}{
\textbf{Status: literature-implied full asymptotic answer.}
For fixed $k\ge2$, $p\ge2$, and $m=\lfloor k/2\rfloor$,
\[
 C(N,p,k)\asymp_{p,k}N^{\max\{0,\,1/2-m/p\}}.
\]
In particular, for every fixed $p>k$, the constants are unbounded and their growth rate is $N^{1/2-\lfloor k/2\rfloor/p}$.}}
\end{center}

\section*{Source questions}
For a $k$-wise independent Rademacher vector $\varepsilon=(\varepsilon_i)_{i=1}^N$, the source defines
\[
C(N,p,k)=\sup_{\substack{\|a\|_2=1\\ \varepsilon\ k\text{-wise independent}}}
 \left(\E\left|\sum_{i=1}^N a_i\varepsilon_i\right|^p\right)^{1/p}.
\]
On source PDF page 2 it asks whether this is bounded for fixed $p,k$ and, if not, for its growth rate.  Page 3 says that understanding $k\ge4$ remains open \cite{PSsource}.  The revised paper restricts the numbered question to $p>k$ and proves the required general upper estimate for even $k$, while again leaving sharpness for $k\ge4$ open \cite{PSrevised}.

\begin{theorem}
Fix $k\ge2$ and $p\ge2$, and put $m=\lfloor k/2\rfloor$.  As $N\to\infty$,
\[
C(N,p,k)\asymp_{p,k}N^{\max\{0,\,1/2-m/p\}}.
\]
Hence $C(N,p,k)$ is bounded in $N$ exactly when $p\le2m$.
\end{theorem}

\section*{Upper bound}
Set $r=2m\le k$.  Any $k$-wise independent Rademacher vector is $r$-wise independent.  If
$f=|\sum_i a_i\varepsilon_i|$, then expansion of the even power and $r$-wise independence give exactly the independent-sign moment:
\[
 \|f\|_r\le C(r)\|a\|_2.
\]
Also $\|f\|_\infty\le\|a\|_1\le\sqrt N\|a\|_2$.  Thus, for $p\ge r$,
\[
 \|f\|_p^p\le \|f\|_\infty^{p-r}\|f\|_r^r,
 \qquad
 \|f\|_p\le C(r)^{r/p}N^{1/2-r/(2p)}\|a\|_2.
\]
This is Proposition 2.1 of \cite{PSrevised}, applied at the largest even order not exceeding $k$.  If $2\le p\le r$, monotonicity of $L_q$ norms gives instead $\|f\|_p\le\|f\|_r\le C(r)\|a\|_2$.  These are the claimed upper bounds in both regimes.

\section*{Lower bound from the maximal endpoint atom}
Let $M(N,k,1/2)$ be the largest possible probability of the all-ones vector under a $k$-wise independent law on $N$ fair bits.  Peled--Yadin--Yehudayoff prove, for every fixed $k$, the sharp asymptotic order
\begin{equation}\label{eq:endpoint}
 M(N,k,1/2)\asymp_k N^{-\lfloor k/2\rfloor}=N^{-m};
\end{equation}
see the odd-$k$ reduction (1.1), Theorem 1.1, and the fixed-$k$ conclusion after Corollary 1.3 on supporting PDF pages 4--7 \cite{PYY}.

Choose a law attaining the maximum (the feasible polytope is compact), write its bits as $X_i\in\{0,1\}$, and set $\varepsilon_i=2X_i-1$.  Then $\varepsilon$ is a $k$-wise independent Rademacher vector.  With $a_i=N^{-1/2}$, on the all-ones event
\[
 \sum_{i=1}^N a_i\varepsilon_i=\sqrt N.
\]
Consequently, by \eqref{eq:endpoint},
\[
 C(N,p,k)^p\ge M(N,k,1/2)N^{p/2}gtrsim_k N^{p/2-m},
\]
and therefore
\[
 C(N,p,k)\gtrsim_{p,k}N^{1/2-m/p}.
\]
This supplies the matching lower bound whenever $p>2m$.  When $p\le2m$, taking $a=e_1$ gives $C(N,p,k)\ge1$, matching the bounded upper estimate.  The theorem follows.

\section*{Proof intuition}
Limited independence fixes every moment only up to the largest even order $2m\le k$.  Interpolating that controlled moment with the deterministic ceiling $\sqrt N$ produces the upper exponent.  The complementary extremal law spends probability about $N^{-m}$ at the most distant Hamming vertex.  Equal coefficients turn this small atom into a value $\sqrt N$; its $p$th-moment contribution is exactly the missing power.  The two estimates meet because both are governed by the same number $m$ of controlled even moments.

\section*{Provenance and scope}
The endpoint estimate \eqref{eq:endpoint} predates the source, but neither \cite{PYY} nor the later Khinchine paper states the combined conclusion above.  Conversely, the revised Pass--Spektor paper gives the matching upper bound but does not cite the endpoint estimate as a lower-bound mechanism.  The result is therefore a direct implication of established literature, not a new construction.

The notation $\asymp_{p,k}$ means two-sided bounds by positive constants depending only on fixed $p,k$.  This fully answers ``bounded or not'' and ``growth rate'' in the usual asymptotic-order sense; it does not determine an exact finite-$N$ optimizer or an exact leading constant for $k\ge4$.

{\raggedright
\begin{thebibliography}{9}
\bibitem{PSsource}
B. Pass and S. Spektor,
\emph{On Khintchine type inequalities for pairwise independent Rademacher random variables},
arXiv:1412.7859v1 (2014), pp.~2--3.

\bibitem{PYY}
R. Peled, A. Yadin, and A. Yehudayoff,
\emph{The Maximal Probability that $k$-wise Independent Bits are All 1},
arXiv:0801.0059v3 (2011), Theorem~1.1 and Corollary~1.3.

\bibitem{PSrevised}
B. Pass and S. Spektor,
\emph{On Khintchine type inequalities for $k$-wise independent Rademacher random variables},
arXiv:1708.08775v1 (2017), Proposition~2.1.
\end{thebibliography}
\par}

\end{document}
